%%%BLOCK id=101-01 ch=15 sec=折射定律·全反射·色散 kind=summary alt=none cont=none y=0.03-0.22
①
% 原稿：左侧一个大括号并列下面三项
\begin{itemize}
\item 折射定律\quad $\dfrac{\sin\theta_1}{\sin\theta_2}=n$。
\item 全反射、光导纤维。
\item 光的色散\quad 材料对不同色光折射率不同\ $n_{\text{红}}$最小；$n_{\text{紫}}$最大，偏折程度最大
\end{itemize}
%%%END
%%%BLOCK id=101-02 ch=15 sec=折射定律和折射率的理解及应用 kind=summary alt=none cont=none y=0.24-0.31
②\ 光的折射。全反射、折射定律和折射率的理解及应用

$n$与介质、光的频率有关。同一色光在不同介质，$v$、$\lambda$ 不同但 $f$ 相同、光路可逆
%%%END
%%%BLOCK id=101-03 ch=15 sec=折射定律和全反射规律的综合应用 kind=summary alt=none cont=none y=0.32-0.40
③\ 折射定律和全反射规律的综合应用
\[ \theta\ge C=\arcsin\frac{1}{n}\qquad v=\frac{c}{n}\qquad t=\frac{L}{v} \]
光路控制、光的色散。
%%%END
%%%BLOCK id=101-04 ch=15 sec=平行玻璃砖·圆柱体·三棱镜 kind=summary alt=none cont=none y=0.42-0.45
④\ 平行玻璃砖\quad 圆柱体、三棱镜
%%%END
%%%BLOCK id=101-05 ch=15 sec=光的波动性·电磁波·相对论 kind=summary alt=14 cont=none y=0.45-0.62
⑤\ 光的波动性、电磁波、相对论
% 原稿：此标题右侧一个大括号括出下面三项
\begin{itemize}
\item 光的\uline{干涉}、衍射、和偏振。
  \begin{itemize}
  \item 加强减弱的区交替。
  \item $\uparrow$光波长略大或与障碍小孔尺寸相近
    % 原稿此项与上一项写在同一行；“↑”箭头指向上方的“衍射”；其右侧有大括号括出以下三项
    \begin{itemize}
    \item 单缝衍射
    \item 圆孔衍射
    \item 圆盘衍射。
    \end{itemize}
  \end{itemize}
\item 电磁波
\item 相对论。
\end{itemize}
%%%END
%%%BLOCK id=101-06 ch=15 sec=考点·光学与电磁波相对论 kind=summary alt=14 cont=none y=0.64-0.80
% 原稿：左缘有一个“L”形标记（疑为序号⑥，未画圈）；页面左缘 y≈0.86 处另有半个圆圈（疑为序号⑦的圈，内无内容）
\unclear{⑥}\ \textbf{考点。}
% 以下三项在原稿中被“考点”右侧的大括号括起
\begin{itemize}
\item 光的干涉现象。
\item 光的衍射和偏振现象
\item 电磁波与相对论
\end{itemize}
%%%END
%%%BLOCK id=102-02 ch=15 sec=测玻璃折射率·插针法 kind=method exp=1 alt=none cont=none y=0.56-1.00
\textbf{测定玻璃折射率\quad 插针法}

\textbf{实验原理与操作}\quad 铺白纸画线、插针、测量
\[ n=\frac{\sin\theta_1}{\sin\theta_2} \]

\textbf{数据处理}
\begin{itemize}
\item 计算法\quad 计算每次折射率，求出平均值$\bar{n}$
\item 图像法
\end{itemize}

%FIG p102_a 0.37 0.722 0.54 0.795
\notefig[0.28]{figures/p102_a.png}

$\sin\theta_1$—$\sin\theta_2$图像\qquad $n=k=\dfrac{\sin\theta_1}{\sin\theta_2}$

\textbf{误差分析}\quad 单位圆法

%FIG p102_b 0.325 0.785 0.495 0.92
\notefig[0.3]{figures/p102_b.png}

\[ n=\frac{EH}{E'H'} \]

\begin{itemize}
\item[\unclear{①}] 根据插针位置画出入射光线，出射光线时有偏差. % 圈号①位于页边被截断，据②③推测
\item[②] 测量时入射角与折射角的测量误差.
\item[③] 作图时，玻璃砖的宽度绘制有误差，画的两边界不平行出现误差.
\end{itemize}

\textbf{注意事项}
\begin{itemize}
\item 玻璃砖要用厚度大的
\item 入射角在$30^\circ\sim60^\circ$之间
\item 大头针要竖直插在白纸上，且距离应尽量大一些
\item 玻璃砖的折射面要画准，不能用玻璃砖界面代替直尺画界线
\end{itemize}
%%%END
%%%BLOCK id=103-01 ch=15 sec=用双缝干涉测光的波长 kind=method exp=1 alt=none cont=none y=0.08-0.38
\textbf{用双缝干涉测光的波长.}

\textbf{实验原理与操作}
\[ \Delta x=\frac{L\lambda}{d} \]
\begin{itemize}
\item $L$为双缝与屏的距离$L$
\item $d$为$S_1,S_2$间距$d$
\item $\Delta x$：相邻两条明条纹间距.
\item $\lambda$：入射光波长
\end{itemize}
\[ \Delta x=\frac{x_2-x_1}{n-1} \]

\textbf{数据处理}

$\lambda=\dfrac{d}{L}\Delta x$\quad 计算波长，重复测量，计算，求出波长的平均值.
%%%END
%%%BLOCK id=103-02 ch=15 sec=用双缝干涉测光的波长 kind=method exp=1 alt=none cont=none y=0.40-0.78
\textbf{误差分析}
\begin{itemize}
\item 双缝到光屏的距离$L$的测量存在误差.
\item 干涉条纹未调到最清晰的程度
\item 分划板中心刻线与干涉条纹不平行，中心刻线没有恰好位于亮条纹中心.
\item 测量多条亮条纹间距间的距离时读数不准确，此问题中条纹数未数清. % CHECK: 原文"间距间的距离"，疑为"间的距离"
\end{itemize}

\textbf{注意事项.}
\begin{itemize}
\item 单缝和双缝要平行，中心位于遮光筒的轴线上.
\item 调节测量头时，应使分划板中心刻线和亮条纹的中心对齐，记下此时手轮上的读数，转动手轮，使分划板中心刻线和另一亮条纹的中心对齐，记下此时手轮上的读数，两次读数之差就表示这两条亮条纹间的距离.
\item 不能直接测$\Delta x$. 要测多条亮条纹间距再计算得到$\Delta x$,
\item 白光的干涉观察到的是彩色条纹，其中白色在中央，红色在最外层
\end{itemize}
%%%END
%%%BLOCK id=104-01 ch=15 sec=薄膜干涉·牛顿环 kind=knowledge alt=none cont=none y=0.05-0.23
\textbf{牛顿环}

\unclear{多个}透镜接触而干涉

%FIG p104_a 0.045 0.124 0.185 0.172
\notefig[0.28]{figures/p104_a.png}
%FIG p104_b 0.285 0.113 0.435 0.226
\notefig[0.3]{figures/p104_b.png}

$\Delta r=2d$ % CHECK: 原稿字形似"Y"(该书写者的r常写成γ/Y形，参见106页r_2-r_1)，按光程差Δr记

\begin{tabular}{@{}c@{\qquad}c@{}}
单色光 & 白光 \\
$\downarrow$ & $\downarrow$ \\
中暗边亮 & 彩 \\
\end{tabular}
\quad 越近越宽. % 原稿"中暗边亮"与"彩"之间有"/"分隔，单色光、白光的箭头分别指向二者
%%%END
%%%BLOCK id=104-02 ch=15 sec=光的衍射·圆盘衍射 kind=knowledge alt=none cont=none y=0.23-0.37
\textbf{圆盘衍射（泊松亮斑）}\quad 光照到半径很小的圆盘上.

%FIG p104_c 0.213 0.2555 0.378 0.352
\notefig[0.28]{figures/p104_c.png}

\begin{itemize}
\item 中心亮斑
\item 阴影外有\unclear{大}不等间距明暗相间的圆环. % CHECK: "大"字疑为"大小"或指示箭头
\end{itemize}
%%%END
%%%BLOCK id=104-03 ch=15 sec=光的衍射·圆孔衍射 kind=knowledge alt=none cont=none y=0.37-0.57
\textbf{圆孔衍射}

%FIG p104_d 0.225 0.400 0.328 0.466
\notefig[0.25]{figures/p104_d.png}

\textbf{单色光}
\begin{itemize}
\item 中央大亮斑\quad 明暗相间的同心圆环
\item 越往外\quad 条亮纹亮度减弱，宽度减小. % CHECK: 原文"条亮纹"，疑为"亮条纹"
\item 亮环或暗环间的距离随圆孔半径增加而减小,
\end{itemize}

\textbf{白光}\quad 中白边彩\quad 中央大且亮的白色亮斑\ 周围是不等距彩色的同心圆环.
%%%END
%%%BLOCK id=104-04 ch=15 sec=光的衍射·单缝衍射 kind=knowledge alt=none cont=none y=0.57-0.80
\textbf{单缝衍射.}

%FIG p104_e 0.12 0.60 0.292 0.692
\notefig[0.3]{figures/p104_e.png}

\textbf{单色光}
\begin{itemize}
\item 中央为亮且宽的条纹，两侧为明暗相间的条纹
\item 越靠近两侧，亮条纹亮度减弱，宽度越小.
\end{itemize}

\textbf{白光、}
\begin{itemize}
\item 中央为亮且宽的白色条纹，两侧为宽度逐渐变窄的彩色条纹.
\item 最靠近中央的是紫光，离中央最远的是红光.
\end{itemize}
%%%END
%%%BLOCK id=104-05 ch=15 sec=光的干涉·双缝干涉 kind=knowledge alt=none cont=none y=0.80-0.92
\textbf{光的干涉}

%FIG p104_f 0.128 0.835 0.31 0.898
\notefig[0.3]{figures/p104_f.png}

条纹宽度相等、等距、亮度相同.
%%%END
%%%BLOCK id=104-06 ch=15 sec=干涉与衍射比较 kind=summary alt=none cont=none y=0.90-1.00
\begin{center}
\begin{tabular}{ll}
\toprule
单缝衍射 & 双缝干涉 \\
\midrule
条纹宽度不等，中央最宽 & 条纹宽度相等. \\
相邻条纹间距不等 & 相邻条纹等间距 \\
中央条纹最亮，两边变暗 & 条纹清晰，亮度基本相等 \\
\bottomrule
\end{tabular}
\end{center}

本质都是光的叠加 % 原稿此句位于表格左侧；页面底边截断，表格最后一行右栏"亮度基本相等"为两行书写
%%%END
%%%BLOCK id=105-01 ch=15 sec=光的偏振 kind=knowledge alt=none cont=none y=0.02-0.24
①\quad \textbf{光的偏振}

\begin{tabular}{@{}l p{0.36\linewidth} p{0.36\linewidth}@{}}
 & 自然光 & 偏振光 \\
来源 & 从普通光源发出的光 & 自然光通过起偏器后的光 \\
 & 在垂直于光的传播方向的平面内，光振动沿任意方向，且沿各方向振动的光强度相同 & 在垂直于光的传播方向的平面内. 光振动沿特定方向，\newline 应用：加偏振滤光片的照相机\unclear{镜头}\newline 液晶显示器、立体电影、\newline 消除车灯眩光 \\
\end{tabular}

%FIG p105_a 0.17 0.118 0.455 0.178
%FIG p105_b 0.47 0.1325 0.545 0.158
\noindent\begin{minipage}[c]{0.6\linewidth}\notefig[1.0]{figures/p105_a.png}\end{minipage}\begin{minipage}[c]{0.2\linewidth}\notefig[1.0]{figures/p105_b.png}\end{minipage}

横波\quad 通过偏振片、反射、折射变成只沿某方向振动的光

$V_{\text{传播}}\perp V_{\text{振动}}$
%%%END
%%%BLOCK id=105-04 ch=15 sec=狭义相对论 kind=knowledge alt=none cont=none y=0.54-0.88
⑤\quad \textbf{相对论.}
\begin{itemize}
\item 狭义相对论
 \begin{itemize}
 \item 狭义相对性原理：在不同的惯性系中一切物理规律相同. % 左页边此行旁有一"L"形短线，疑为层级连线
 \item 光速不变原理\quad 真空中的光速在不同惯性参考系中都相同，与光源、观测者间相对运动无关.
 \end{itemize}
\end{itemize}
\[ \left\{\begin{array}{l} m=\dfrac{m_0}{\sqrt{1-\left(\dfrac{v}{c}\right)^{2}}} \\[3ex] E=mc^{2} \end{array}\right. \]

高速下\quad 尺缩\quad 钟慢\quad 质增 % 其后有涂改带遮盖，字迹不可辨
%%%END
%%%BLOCK id=106-01 ch=15 sec=双缝干涉·亮暗条纹判断 kind=knowledge alt=none cont=none y=0.03-0.22
\textbf{判断双缝干涉亮暗条纹的方法}

%FIG p106_a 0.115 0.045 0.335 0.15
\notefig[0.4]{figures/p106_a.png}

对于两步调一致的光源 % CHECK: 原文"两步调"，疑漏"个"

$r_2-r_1=k\lambda\ (k=0,1,2)$\quad 光屏上出现亮条纹 % 原稿"2"后仅一点，疑为省略号

$r_2-r_1=(2k+1)\dfrac{\lambda}{2}\ (k=0,1,2\ldots)$\quad 光屏上出现暗条纹,

对于两步调相反的光源，产生亮暗纹与上面的正好相反
%%%END
%%%BLOCK id=106-02 ch=15 sec=薄膜干涉 kind=knowledge alt=none cont=none y=0.23-0.44
\textbf{薄膜干涉}

%FIG p106_b 0.02 0.232 0.49 0.405
\notefig[0.55]{figures/p106_b.png}

$\Delta r=n\lambda$\quad $n=1,2,3$\quad 亮条纹

$\Delta r=(2n+1)\dfrac{\lambda}{2}$\quad $(n=0,1,2\ldots)$\quad 暗条纹

由重力作用成上薄下厚的楔形
%%%END
%%%BLOCK id=106-03 ch=15 sec=干涉法检查平面 kind=knowledge alt=none cont=none y=0.44-0.58
\textbf{干涉法检查平面.}

%FIG p106_c 0.125 0.472 0.495 0.573
\notefig[0.5]{figures/p106_c.png}

\begin{itemize}
\item 若平面平整，则观察到平行等间距的明暗相间条纹 % 原稿"平面"二字连写成一个字形(似"辅")，据文意读作"平面"
\item 若平面不平整，则干涉条纹发生弯曲
\end{itemize}
%%%END
%%%BLOCK id=111-02 ch=15 sec=折射定律 kind=knowledge alt=none cont=none y=0.125-0.23
\textbf{光}\\[0.8ex]
折射律　$n=\dfrac{\sin i}{\sin r}$
%FIG p111_b 0.31 0.139 0.44 0.205
\notefig[0.25]{figures/p111_b.png}
不是
%FIG p111_c 0.545 0.143 0.655 0.222
\notefig[0.25]{figures/p111_c.png}
\begin{align*}
n_1&=\dfrac{\sin i_1}{\sin r_1}\\
n_2&=\dfrac{\sin i_2}{\sin r_2}\qquad V=\dfrac{c}{n}
\end{align*}
%%%END
%%%BLOCK id=111-03 ch=15 sec=全反射 kind=knowledge alt=none cont=none y=0.26-0.41
% 下一行为写在“光密→光疏”上方的小字批注
水/玻璃　空气　\unclear{空气棱镜}$C$相除$n$\\
全反射　（光密$\to$光疏　不折射　只反射的现象）
\[
\sin C=\dfrac{1}{n}
\]
%FIG p111_d 0.405 0.32 0.515 0.388
\notefig[0.25]{figures/p111_d.png}
\[
n=\dfrac{\sin 90^\circ}{\sin C}=\dfrac{1}{\sin C}
\]
\begin{center}
\begin{tabular}{cccccc}
$n$ & $\dfrac{2\sqrt{3}}{3}$ & $\sqrt{2}$ & $2$ & $1.5$ & $\sqrt{3}$\\[1.2ex]
$C$ & $60^\circ$ & $45^\circ$ & $30^\circ$ & $42^\circ$ & $35^\circ$
\end{tabular}
\end{center}
%%%END
%%%BLOCK id=111-04 ch=15 sec=全反射 kind=method alt=none cont=none y=0.41-0.52
\textbf{Step 1}
\begin{itemize}
\item 分界面　$n=\dfrac{\sin i}{\sin r}$
\item 临界　$\sin C=\dfrac{1}{n}$
\item 反射　$\alpha=\beta$
\end{itemize}
画光路图判定是否全反射　画反射光线'\\[0.8ex]
\textbf{Step 2}　几何角度关联
\begin{itemize}
\item 三角函数、外角、正弦定理、余弦定理　\unclear{勾股}
\item 用量角器量角度
\end{itemize}
$\cos2\theta=2\cos^2\theta-1$
%%%END
%%%BLOCK id=111-05 ch=15 sec=折射率 kind=formula alt=none cont=none y=0.51-0.67
\[
n=\dfrac{n_{\text{介质1}}}{n_{\text{介质2}}}=\dfrac{\sin\theta_{\text{大}}}{\sin\theta_{\text{小}}}=\dfrac{\lambda_{\text{真空}}}{\lambda_{\text{介质水}}}=\dfrac{c}{v}=\dfrac{1}{\sin C}>1
\qquad n_{\text{介质}}>1\quad n_{\text{空气}}=1
\]
%FIG p111_e 0.15 0.563 0.255 0.625
\notefig[0.25]{figures/p111_e.png}
\[
n_1\sin\theta_1=n_2\sin\theta_2\quad\text{（完整版公式）}
\]
$n$与介质和光频率有关\\[1ex]
光在　空气中传播时间　$t=\dfrac{s}{c}$　　介质中传播　$t=\dfrac{s}{v}$
%%%END
%%%BLOCK id=111-06 ch=15 sec=视深 kind=knowledge alt=none cont=none y=0.65-0.83
水面垂直视深度　$h=\dfrac{h_{\text{实际}}}{n}$　（原高度除以折射率）\\[0.8ex]
球面垂直视深度　$h'=h$（不变）
%FIG p111_f 0.578 0.672 0.80 0.80
\notefig[0.4]{figures/p111_f.png}
$\theta$很小时　$\tan\theta\approx\sin\theta$
\begin{align*}
\tan\theta_1&=\dfrac{L}{h'}\\
\tan\theta_2&=\dfrac{L}{h}
\end{align*}
% CHECK: tanθ1=L/h'中的分母字迹像n'，按物理应为h'
\[
n=\dfrac{\sin\theta_1}{\sin\theta_2}=\dfrac{\tan\theta_1}{\tan\theta_2}=\dfrac{h}{h'}
\]
%%%END
%%%BLOCK id=111-07 ch=15 sec=棱镜与玻璃砖 kind=knowledge alt=none cont=none y=0.80-0.97
三棱镜　两面均色散，偏折向厚边\\
圆镜　半径为法线　同心圆用$2$个\unclear{根}正弦定理\\
方砖玻璃　出入平行，两面都不会全反射，两面都色散　　$s=\dfrac{d}{\cos\alpha}$　　$t=\dfrac{s}{v}=\dfrac{dn}{\cos\alpha}$ % CHECK: 原稿t=dn/cosα，缺光速c（应为dn/(c cosα)），照原样转录
%FIG p111_g 0.41 0.895 0.537 0.972
\notefig[0.25]{figures/p111_g.png}
侧移量　$x\propto\theta,\ d,\ n.$　\fbox{\unclear{$\uparrow d+\uparrow n$}}　同大\unclear{小}
%%%END
%%%BLOCK id=112-01 ch=15 sec=折射·全反射 kind=note alt=none cont=prev y=0.04-0.09
% DUP?: 页眉区（本页第一条横线之上）的字迹，其中“t d + En 同大小”与 p111 页底最后一行“(td+En) 同大小”相同，可能是相邻页字迹露出；若与 p111 重复可删除
\unclear{似保?}\quad \unclear{$t\,d\,+E_n$} 同大小
%%%END
%%%BLOCK id=112-02 ch=15 sec=光的衍射 kind=knowledge alt=none cont=none y=0.10-0.18
\textbf{光的衍射}．波长大衍射强，绕过障碍，$\lambda>$小孔尺\unclear{寸}时/障碍物尺寸时 明显衍射

光不沿直线传播了
%%%END
%%%BLOCK id=112-03 ch=15 sec=光的干涉·双缝 kind=knowledge alt=none cont=none y=0.14-0.42
\textbf{光的干涉}：同频(色)光、相位差不变、振动方向相同、只有同频率才能干涉

\struck{同\unclear{频}光衍射的叠加}\quad 同频光叠加

%FIG p112_a 0.213 0.168 0.39 0.267
\notefig[0.3]{figures/p112_a.png}

亮纹\quad $\Delta S=\dfrac{\lambda}{2}(2n)$\quad $n=0$ 第0亮纹 ...

暗纹\quad $\Delta S=\dfrac{\lambda}{2}(2n+1)$\quad \fbox{对一个\unclear{已知角}用余弦定理}

\[\Delta x=\frac{L'}{n-1}=\frac{L\lambda}{d}\] % CHECK: 分子疑为 L' 或 L 加一撇，原稿不清
$n$为$n$组条纹间距\quad 考察单位换算，数量级 % CHECK: “n为n组条纹间距”原文如此

双缝向下，$P$也向下（保同时性）

单色光\quad 红黑等距\quad 中亮边暗

复色光\quad 彩色等距\quad 中白边彩\quad 红光最外

绿光$+$蓝光$\Rightarrow$独立的蓝光条纹、绿光条纹
%%%END
%%%BLOCK id=112-04 ch=15 sec=光的干涉·薄膜 kind=knowledge alt=none cont=none y=0.42-0.55
\textbf{液体膜}(水/肥皂泡)

条纹间距\quad $x=\dfrac{\lambda}{2\tan\theta}=\dfrac{\lambda}{2\sin\theta}=\dfrac{\lambda}{2\theta}$\quad $\theta$很小\quad 液体膜顶角

%FIG p112_b 0.18 0.482 0.285 0.545
\notefig[0.25]{figures/p112_b.png}

前膜反射光和后膜反射光相干涉

光在人的一侧
%%%END
%%%BLOCK id=112-05 ch=15 sec=光的干涉·薄膜 kind=knowledge alt=none cont=none y=0.55-0.71
\textbf{气体膜}

% 图p112_c内含标注：窄 宽 窄（上方）；各处θ不同（下方）
%FIG p112_c 0.10 0.555 0.225 0.644
\notefig[0.3]{figures/p112_c.png}

% 图p112_d：圆环（牛顿环俯视），中心标“暗”
%FIG p112_d 0.17 0.643 0.27 0.71
\notefig[0.25]{figures/p112_d.png}

$x=\dfrac{\lambda}{2\tan\theta}$

$\theta$大条纹变密，$\theta$为气体膜顶角\quad 曲率

中间暗，往外为亮暗相间

% 图p112_e：楔形气体膜（带上下箭头）
%FIG p112_e 0.67 0.595 0.79 0.655
\notefig[0.25]{figures/p112_e.png}

% 图p112_f内含标注：人俯视条纹
%FIG p112_f 0.85 0.585 0.97 0.67
\notefig[0.25]{figures/p112_f.png}

\begin{tabular}{@{}cc@{\qquad}}
左凹 & 右凸 \\
同亮 & 同厚
\end{tabular}
%%%END
%%%BLOCK id=112-06 ch=15 sec=光的干涉·薄膜 kind=knowledge alt=none cont=none y=0.68-0.79
\textbf{固体膜}

\begin{itemize}
\item 增透膜(减反膜)\quad 增强透射光，减少反射光\\
  $2d=\dfrac{1}{2}\lambda\,(2n+1)$\qquad $d=\dfrac{1}{4}\lambda\cdot(2n+1)$
\item 增反膜(减透膜)\quad 增强反射光，减少透射光\\
  $2d=\dfrac{1}{2}\lambda\,(2n)$\qquad $d=\dfrac{1}{4}\lambda\cdot 2n$,
\end{itemize}

%FIG p112_g 0.555 0.72 0.645 0.785
\notefig[0.25]{figures/p112_g.png}
%%%END
%%%BLOCK id=112-07 ch=15 sec=光的衍射 kind=knowledge alt=none cont=none y=0.80-0.93
\textbf{单缝衍射、小孔衍射}，中间宽，两边暗．中宽边窄．中白边彩，红光最外

$\left\{\begin{tabular}{@{}l@{}}单缝越窄\\波长越长\end{tabular}\right.$ $\to$衍射越明显$\to$条纹间距大$\to$亮条纹范围大$\to$亮条纹亮度\unclear{低}

\textbf{圆盘衍射}(泊松亮斑)\quad 中窄边宽，中亮\quad 亮暗相间
%%%END
%%%BLOCK id=112-08 ch=15 sec=光的偏振 kind=knowledge alt=none cont=next y=0.95-1.00
% 本页底部被拍摄边缘截断，最后一行只露出上半部分
\textbf{偏振}：横波特有\quad 只通过平行光栅振动方向的光

偏振片可滤去反射光\quad 当偏振方向与反射光的偏振方向垂直时\unclear{…}(\unclear{横}波)\unclear{光…}

横波\quad 振动方向$\perp$传播方向

纵波\quad 振动方向$\parallel$传播方向
%%%END
%%%BLOCK id=113-03 ch=15 sec=折射·全反射 kind=knowledge alt=none cont=none y=0.22-0.35
\begin{enumerate}
\item 彩虹为光的色散/折射
\item \unclear{鸟}看鱼偏高、鱼看鸟偏高
\item 水中唱歌\quad 从空气到水中\ 声波波长$\uparrow$，电磁波波长短
\item 海\unclear{市蜃}楼、车尾灯光带(全反射)
\end{enumerate}
%%%END
%%%BLOCK id=113-04 ch=15 sec=折射·全反射 kind=problem alt=none cont=none y=0.35-0.45
% 原稿第5题：只有图与解答，没有写出题干
\begin{problem}[5]
%FIG p113_a 0.09 0.352 0.205 0.43
\notefig[0.25]{figures/p113_a.png}
\begin{solution}
\begin{align*}
&\sin\alpha=\frac{1}{n_1}=\frac{r_1}{\sqrt{r_1^2+d^2}}\quad\therefore\ r_1=\frac{d}{\sqrt{n_1^2-1}}\qquad
\sin\beta=\frac{1}{n_2}=\frac{r_2}{\sqrt{r_2^2+d^2}}\qquad r_2=\frac{d}{\sqrt{n_2^2-1}}\\
&S=\pi r_2^2-\pi r_1^2=\pi d^2\left(\frac{1}{n_2^2-1}-\frac{1}{n_1^2-1}\right)
\end{align*}
\end{solution}
\end{problem}
%%%END
%%%BLOCK id=113-05 ch=15 sec=折射·全反射 kind=derivation alt=none cont=none y=0.45-0.66
\textbf{6．光导纤维．}

%FIG p113_b 0.055 0.478 0.25 0.567
\notefig[0.35]{figures/p113_b.png}

$n_1>n_2$
\[n_1\sin\theta_3=n_2\sin\frac{\pi}{2}\qquad\therefore\ \sin\theta_3=\cos\theta_2\qquad\therefore\ \theta_3+\theta_2=\frac{\pi}{2}\]
\[\left\{\begin{array}{l}n_1=\dfrac{\sin\theta_1}{\sin\theta_2}\\[2mm]\sin^2\theta_2+\cos^2\theta_2=1\quad\therefore\ \sin\theta_1=\sqrt{n_1^2-n_2^2}\end{array}\right.\]

路程\ $x=\dfrac{L}{\cos\theta_2}$\qquad $v=\dfrac{c}{n_1}$\qquad $t=\dfrac{x}{v}=\dfrac{n_1^2}{n_2^2}\,\dfrac{L}{c}$ % CHECK: 分母似为 n_2^2，但由前式 cosθ2=n2/n1 推得应为 n_1^2/n_2 · L/c

% 以下两式在图的左下方；第一式整体被划去
\[\text{\struck{$\dfrac{1}{\sin\theta_3}=n=\dfrac{\sin\theta_1}{\sin\theta_2}$}}\]
\[\frac{n_1}{n_2}=\frac{1}{\sin\theta_3}\]
%%%END
%%%BLOCK id=113-06 ch=15 sec=折射·全反射 kind=problem alt=none cont=none y=0.66-0.92
% 原稿第7题：只有图与解答，没有写出题干
\begin{problem}[7]
% 图p113_c内含标注：中点、中线、H、ΔH、V0、水位上升速度V0、D（底面宽）、Δd、d、α、β
%FIG p113_c 0.0 0.675 0.262 0.832
\notefig[0.4]{figures/p113_c.png}
\begin{solution}
① 求折射率
\begin{align*}
&\sin\alpha=\frac{D}{\sqrt{D^2+H^2}}\qquad \sin\beta=\frac{d-\frac{D}{2}}{\sqrt{\left(d-\frac{D}{2}\right)^2+\left(\frac{H}{2}\right)^2}}\\
&n=\frac{\sin\alpha}{\sin\beta}=\frac{D}{2d-D}\sqrt{\frac{(2d-D)^2+H^2}{D^2+H^2}}
\end{align*}
② 求光源与水面速度关系
% 下式（①式）下方有红色波浪线
\begin{align*}
&\Delta d=\Delta H(\tan\alpha-\tan\beta)\quad\text{①}\\
&\tan\alpha=\frac{D}{H}\qquad \tan\beta=\frac{d-\frac{D}{2}}{\frac{H}{2}}\quad\nearrow\\
&\Delta d=\Delta H\,\frac{(2H-2d)}{H}\quad\text{②}\\
&\frac{\Delta d}{\Delta t}=\frac{\Delta H}{\Delta t}\,\frac{(2\unclear{H}-2d)}{H}\quad\text{③}\\
&\therefore V=\frac{2(D-d)}{H}V_0\quad\text{④}
\end{align*}
% CHECK: ②③式分子中的“2H”疑应为“2D”（由 tanα-tanβ=(2D-2d)/H 及④式2(D-d)/H）；③式中该字母似为D写在H之上。tanβ 行末有一弯曲箭头指向上方①式
\end{solution}
\end{problem}
%%%END
%%%BLOCK id=195-04 ch=15 sec=双缝干涉·条纹间距与波长 kind=knowledge exp=1 alt=none cont=none y=0.88-0.94
相邻$n$个亮条纹中心距离为$\Delta x$，则条纹间距为\ $\dfrac{\Delta x}{n-1}=\dfrac{\unclear{L\lambda}}{d}$\qquad $\lambda=\dfrac{\Delta x\,d}{3L}$
%%%END
%%%BLOCK id=196-03 ch=15 sec=折射·视深与色散干涉现象 kind=note alt=none cont=none y=0.40-0.44
鸟看鱼\ 鱼看鸟\ 都偏高\qquad \fbox{彩虹\textsuperscript{折射}\ \ 皂泡\textsuperscript{干涉}}\ $\in$\ 色\struck{彩}散
% 原稿“折射”“干涉”为写在“彩虹”“皂泡”右上方的小字；“彩虹…干涉”四项被一个椭圆圈起；“色散”中间一字被涂改
%%%END
%%%BLOCK id=226-01 ch=15 sec=折射定律·折射率 kind=note alt=none cont=none y=0.03-0.10
\textbf{1.}
%FIG p226_a 0.10 0.02 0.25 0.09
\notefig[0.3]{figures/p226_a.png}

\unclear{NO} $n=\dfrac{\sin 2\theta}{\sin\theta}$ % CHECK: 公式左上角的小字"NO"含义不明
%%%END
%%%BLOCK id=228-04 ch=15 sec=反射·两镜多次反射 kind=note alt=none cont=none y=0.62-0.70
\textbf{3.}
%FIG p228_d 0.07 0.625 0.155 0.695
\notefig[0.2]{figures/p228_d.png}

\red{两\unclear{镜}区 多次反射问题}
%%%END
%%%BLOCK id=228-05 ch=15 sec=折射·半圆形玻璃砖 kind=note alt=none cont=none y=0.62-0.73
%FIG p228_e 0.47 0.62 0.595 0.725
\notefig[0.25]{figures/p228_e.png}

出入平行 必过最高点

$\dfrac{\sin i}{\sin r}=n$
%%%END
%%%BLOCK id=232-06 ch=15 sec=光传播时间·空气另算 kind=note alt=11 cont=none y=0.445-0.525
\textbf{6、}
%FIG p232_d 0.095 0.438 0.164 0.52
\notefig[0.2]{figures/p232_d.png}
\begin{keybox}
空气时间另算 // 无\unclear{场}区匀速阶段
\end{keybox}
空气透镜 \quad $t=\text{\unclear{$\bullet$}}\dfrac{s}{c}$ \quad $n=1$ % CHECK: 分式前有一个辨认不清的小符号（像带竖线的圆点）
%%%END
